% Response-to-reviewers template for the revision round.
% Pre-filled with answers to the questions reviewers are most likely to raise;
% replace "Comment R?.?" with the reviewers' actual wording and keep only what applies.
\documentclass[11pt]{article}
\usepackage[margin=2.2cm]{geometry}
\usepackage{mathptmx}
\usepackage{xcolor}
\usepackage[hidelinks]{hyperref}
\newcommand{\rc}[1]{\par\medskip\noindent\textbf{Comment #1.}\ }
\newcommand{\resp}{\par\noindent\textcolor{blue!60!black}{\textbf{Response.}}\ }
\newcommand{\chg}{\par\noindent\textit{Changes in the manuscript:}\ }
\begin{document}
\begin{center}
{\large\bfseries Response to the Reviewers}\\[2pt]
Manuscript ID: [ACCESS-2026-XXXXX] --- ``Excitation-Coherent Telemetry Analysis for Predictive
OBD-II Connector Wear Detection: A Physics-Informed Study on Real In-Vehicle Network Data''
\end{center}

We thank the Editor and the Reviewers for their careful reading and constructive comments. Below,
each comment is reproduced in bold, followed by our response and the corresponding changes.
Changes in the revised manuscript are highlighted in blue.

\section*{Reviewer 1}

\rc{R1.1 (realism of the faults)} \emph{The faults are simulated; the detector may only learn the simulator.}
\resp We agree that this is the central question and address it in four ways. (i) All healthy
telemetry, the excitation that drives the faults and the usage histories are real; only the
connector fault is modelled, and the detector never sees model internals. (ii) Sixteen parameter
perturbations (Section~VI-H, Fig.~9; Supplement S6): on full-bus data the AUROC stays within
0.84--0.95 and severe-stage AUROC $\ge 0.97$. (iii) A \emph{structurally different} fault process
(bursty Markov chatter with an energy-based brown-out criterion; Section~VI-I, Supplement S8):
trained only on the nominal model, ECTA reaches 0.856 on full-bus data (moderate 0.848, severe
0.999), 0.779 on the polling corpus and 0.813 on VED. (iv) The analytical observability law
(Section~VI-C), which involves no classifier, predicts the measured detectability (Spearman 0.93).
\chg [Section/page references.]

\rc{R1.2 (stronger deep baselines)} \emph{Compare with a Transformer.}
\resp A Transformer encoder baseline is now reported for all four corpora (Section~VI-I,
Supplement S10). On unseen vehicles it reaches 0.641 (full-bus), 0.665 (sampling), 0.665 (polling)
and 0.590 (VED), against 0.880, 0.675, 0.723 and 0.775 for ECTA.

\rc{R1.3 (statistical significance)}
\resp Paired vehicle-level bootstrap confidence intervals and p-values for ECTA against every
baseline are given in Supplement S11. On the full-bus corpus (six vehicles) and VED (381 vehicles)
ECTA's advantage over every non-ECTA baseline is significant at $p\le0.001$ (the resolution of
1\,000 resamples); the three-vehicle corpora support no formal test and are reported as indicative.

\rc{R1.4 (hyperparameters)}
\resp Hyperparameters were fixed before evaluation and are identical across corpora. Across seven
gradient-boosting configurations (Supplement S9) the overall AUROC varies by at most 0.004
(full-bus), 0.007 (polling) and 0.015 (VED).

\section*{Reviewer 2}

\rc{R2.1 (practical false-alarm rate)}
\resp Supplement S12 reports false alarms per hour of real healthy driving at fixed detection
rates. A full-bus tool detecting 90\% of severely worn connectors raises no false alarm in the
held-out healthy data (0.39 per hour at 95\% detection); low-rate polling tools need multi-trip
pooling (Fig.~7), after which severe wear is detected per vehicle with AUROC 0.999.

\rc{R2.2 (window length)}
\resp Supplement S13 reports 10-, 30- and 60-s windows (AUROC 0.824, 0.879, 0.911). Detection
improves with window length as the physics predicts; the 30-s window is a latency compromise, not
the most favourable choice.

\rc{R2.3 (validation on naturally worn connectors)}
\resp We agree that this is the most valuable next step and state it as the primary limitation.
[If bench data have been collected: describe specimens, protocol and results here.]

\section*{Reviewer 3}

\rc{R3.1 [comment]}
\resp [response]
\chg [changes]

\end{document}
